\chapter{BACKGROUND AND RELATED WORK}

\section{BRIDGE WEIGH-IN-MOTION SYSTEMS}

Bridge Weigh-In-Motion (B-WIM) is a non-intrusive traffic monitoring technology
that estimates axle loads and gross vehicle weights of vehicles in motion. The
principle, first described by Moses \cite{Moses1979}, is to instrument a bridge
with strain gauges and use the measured dynamic response together with an
influence line algorithm to infer the forces applied by passing axles.

A typical B-WIM installation consists of strain gauges attached to the soffit of
a bridge girder, a data acquisition system sampling at 100--1000\,Hz, and signal
processing software. As a vehicle crosses, each axle generates a characteristic
strain pulse. The timing and relative amplitude of these pulses depend on the
vehicle's axle spacing, speed, and individual axle weights. Once the axle
positions in time are known, the system solves a system of equations to recover
the individual axle weights \cite{Soberl2025}.

Modern B-WIM systems can achieve gross vehicle weight accuracy within a few
percent under favourable conditions, but performance degrades when signals are
noisy, vehicle speeds vary widely, or multiple vehicles are present on the bridge
simultaneously. The initial step of reliably detecting individual axle events is
therefore critical: a single missed or spurious detection directly corrupts the
weight estimation for that vehicle.

\section{TRADITIONAL AXLE DETECTION}

Classical axle detection treats the smoothed strain signal as a quasi-static
response and identifies axle positions as local maxima. Typical steps are:

\begin{enumerate}
    \item \textbf{Low-pass filtering} to suppress high-frequency vibration noise.
    \item \textbf{Savitzky-Golay smoothing} \cite{SavitzkyGolay1964}, which fits
          a polynomial to a sliding window. This approach preserves peak shapes
          and curvature better than a simple moving average, making subsequent
          peak-finding more reliable.
    \item \textbf{Peak detection} based on manually tuned thresholds for minimum
          peak height, minimum prominence, and minimum inter-peak distance.
\end{enumerate}

These methods are fast and interpretable, but require per-installation parameter
tuning. They are sensitive to signal-to-noise ratio and may miss closely-spaced
axles (e.g., tandem or tridem axle groups) or generate false positives at the
start and end of a crossing event.

\section{DEEP LEARNING FOR TIME SERIES}

Deep learning has achieved state-of-the-art results across many time-series
tasks, including classification, segmentation, anomaly detection, and event
localisation. Two families of architectures are directly relevant to this work.

\subsection{1D Convolutional Neural Networks and U-Net}

One-dimensional CNNs apply learned filters along the temporal axis to extract
local features from a sequence. Stacking multiple convolutional layers with
increasing receptive fields allows the network to capture patterns at multiple
temporal scales.

The U-Net architecture \cite{Ronneberger2015}, originally proposed for biomedical
\emph{image} segmentation, adapts naturally to 1D sequence labelling. An encoder
path successively down-samples the sequence while increasing feature depth; a
symmetric decoder path up-samples back to the original resolution, with skip
connections that transfer fine-grained spatial detail from encoder to decoder
layers. The final output is a per-sample probability map, making U-Net directly
applicable to dense binary labelling (axle / non-axle).

\subsection{Temporal Convolutional Networks}

Temporal Convolutional Networks (TCNs) \cite{Bai2018} replace recurrent
connections with stacks of \emph{dilated} 1D convolutions. Dilation inserts gaps
between filter taps, exponentially increasing the effective receptive field with
network depth without increasing the number of parameters. For a TCN with kernel
size $k$ and $N$ dilated residual blocks (dilation doubling each block), the
total receptive field is:

\begin{equation}
    RF = 1 + 2(k-1)\bigl(2^{N}-1\bigr).
    \label{eq:rf}
\end{equation}

With $k=3$ and $N=9$, Equation~\ref{eq:rf} gives $RF=2{,}045$ samples, which
comfortably covers the entire 1,300-sample signal in a single forward pass.

Unlike recurrent networks, TCNs process all time steps in parallel and are
therefore stable to optimise and fast to train. Bai et al.\ \cite{Bai2018} showed
on a suite of sequence modelling benchmarks that TCNs match or exceed the
performance of LSTMs and GRUs at lower computational cost.

\section{CLASS IMBALANCE IN BINARY SEQUENCE LABELLING}

A critical challenge in axle detection is extreme class imbalance. In the dataset
used in this thesis, only approximately 0.35\,\% of all signal samples correspond
to an axle event, yielding a negative-to-positive sample ratio of 287:1. If this
imbalance is not addressed, a model that always predicts ``no axle'' achieves
99.65\,\% sample-level accuracy while being completely useless.

A standard remedy is to weight the positive class in the loss function. With
binary cross-entropy (BCE) loss and a scalar weight $w^{+}$, the loss becomes:

\begin{equation}
    \mathcal{L} = -\frac{1}{N}\sum_{i=1}^{N}
        \bigl[ w^{+}\, y_i \log \hat{p}_i
              + (1-y_i) \log(1-\hat{p}_i) \bigr],
    \label{eq:bce}
\end{equation}

where $y_i \in \{0,1\}$ is the ground-truth label and $\hat{p}_i$ is the
predicted probability. Setting $w^{+} = 287$ brings the effective contribution of
positive and negative samples into balance during training.

\section{RELATED WORK}

\v{S}oberl et al.\ \cite{Soberl2025} combine strain-based B-WIM heuristics with
computer vision (traffic camera images) to verify axle configurations and flag
potentially incorrect GVW estimates. On a dataset of over 30,000 HGV records
collected from the same B-WIM installation used in this thesis, their combined
approach improves GVW reliability from 96.7\,\% to 99.89\,\%. Their work
highlights the difficulty of axle configuration inference from strain signals
alone and motivates the use of data-driven methods for axle detection.

To the best of the author's knowledge, the direct application of 1D U-Net or TCN
architectures to per-sample axle detection from B-WIM strain signals has not
previously been reported at the bachelor level, making this thesis an original
empirical contribution to the field.